% Section 5: Results and Discussion
% Ready to paste into Springer LNCS manuscript
% All values sourced from ablation_results_per_seed.csv, paired_comparisons.csv,
% threshold_sweep_per_seed.csv, threshold_sweep_trend.csv
% Convention: positive class = APPROVAL (target_default == 0)

\section{Results and Discussion}

All experiments use a synthetic credit-risk dataset ($N{=}5{,}000$) with a
70/30 train--test split ($N_{\text{test}}{=}1{,}500$). Results are reported as
mean $\pm$ standard deviation over 10 independent seeds (42--51), varying
the train--test partition and the Random Forest random state while holding the
dataset fixed. The positive class is \emph{loan approval}
($\texttt{target\_default}{=}0$); the default prevalence in the test split is
$11.57\% \pm 0.78\%$.

\subsection{Overall Performance}

Table~\ref{tab:overall} presents the full metric profile of the proposed
multi-agent framework (C1) alongside the unconstrained Random Forest baseline
(C4) and an always-approve constant predictor (C0a).

\begin{table}[htbp]
\centering
\caption{Overall system performance (mean $\pm$ std, $n{=}10$ seeds).
Positive class = approval. Specificity = share of true defaulters correctly
rejected.}
\label{tab:overall}
\setlength{\tabcolsep}{3pt}
\begin{tabular}{lccc}
\hline
\textbf{Metric} & \textbf{C1: Full Framework} & \textbf{C4: RF Baseline} & \textbf{C0a: Always-Approve} \\
\hline
Accuracy            & $0.8778 \pm 0.0065$ & $\mathbf{0.9021 \pm 0.0063}$ & $0.8843 \pm 0.0078$ \\
Precision           & $\mathbf{0.9262 \pm 0.0065}$ & $0.9204 \pm 0.0058$ & $0.8843 \pm 0.0078$ \\
Recall              & $0.9364 \pm 0.0042$ & $\mathbf{0.9734 \pm 0.0029}$ & $\mathbf{1.0000 \pm 0.0000}$ \\
F1 Score            & $0.9313 \pm 0.0037$ & $\mathbf{0.9462 \pm 0.0037}$ & $0.9386 \pm 0.0044$ \\
ROC-AUC             & $\mathbf{0.8697 \pm 0.0113}$ & $0.8587 \pm 0.0154$ & $0.5000 \pm 0.0000$ \\
Balanced Accuracy   & $\mathbf{0.6832 \pm 0.0202}$ & $0.6652 \pm 0.0113$ & $0.5000 \pm 0.0000$ \\
Specificity         & $\mathbf{0.4300 \pm 0.0392}$ & $0.3569 \pm 0.0202$ & $0.0000 \pm 0.0000$ \\
MCC                 & $0.3802 \pm 0.0395$ & $\mathbf{0.4289 \pm 0.0298}$ & $0.0000 \pm 0.0000$ \\
$\Delta_{\text{DP}}$ & $\mathbf{0.0400 \pm 0.0184}$ & $0.0619 \pm 0.0123$ & $0.0000 \pm 0.0000$ \\
$\Delta_{\text{EO}}$ & $\mathbf{0.1506 \pm 0.0851}$ & $0.2439 \pm 0.0778$ & $0.0000 \pm 0.0000$ \\
Policy Violations   & $\mathbf{20.3 \pm 4.3}$ & $78.0 \pm 9.0$ & --- \\
\hline
\end{tabular}
\end{table}

Table~\ref{tab:overall} reveals a structured trade-off between predictive
accuracy and governance compliance. The Random Forest baseline (C4) attains
higher accuracy ($0.9021$ vs.\ $0.8778$), F1 ($0.9462$ vs.\ $0.9313$), and
MCC ($0.4289$ vs.\ $0.3802$) than the proposed framework (C1). This gap
arises because the policy and fairness enforcement layers in C1 reject or
reroute applications that the unconstrained classifier would approve,
reducing recall from $0.9734$ to $0.9364$.

The always-approve constant predictor (C0a) achieves accuracy $0.8843$ and
F1 $0.9386$, both exceeding C1. This is a direct consequence of the class
imbalance: with $88.43\%$ of test instances belonging to the approval class,
a trivial strategy that approves every application obtains high accuracy and
F1 by construction, while providing zero discriminative ability
(ROC-AUC~$= 0.500$, specificity~$= 0.000$, MCC~$= 0.000$).

The framework recovers ground on metrics that penalise indiscriminate
approval. Specificity---the proportion of true defaulters correctly
rejected---rises from $0.3569$ (C4) to $0.4300$ (C1), and balanced accuracy
from $0.6652$ to $0.6832$. These gains reflect the policy agent's capacity
to block high-risk applicants that the classifier alone would approve.
On the fairness dimensions, $\Delta_{\text{DP}}$ decreases from $0.0619$
(C4) to $0.0400$ (C1), and $\Delta_{\text{EO}}$ from $0.2439$ to $0.1506$.
Policy violations fall from $78.0$ to $20.3$, a $74.0\%$ reduction.

\subsection{Ablation Study}

To isolate the contribution of each agent, four configurations were
evaluated. C1 denotes the full framework; C2 disables the fairness
(counterfactual) agent while retaining the policy agent; C3 disables the
policy agent while retaining the fairness agent; C4 removes both agents,
leaving the Random Forest classifier alone. Table~\ref{tab:ablation}
summarises the results.

\begin{table}[htbp]
\centering
\caption{Controlled ablation study (mean $\pm$ std, $n{=}10$ seeds,
$N_{\text{test}}{=}1{,}500$, $\varepsilon{=}6.0$).}
\label{tab:ablation}
\setlength{\tabcolsep}{2.5pt}
\footnotesize
\begin{tabular}{lcccccccc}
\hline
\textbf{Config} & \textbf{Acc.} & \textbf{F1} & \textbf{Spec.} & \textbf{Bal.\ Acc.} & \textbf{$\Delta_{\text{DP}}$} & \textbf{$\Delta_{\text{EO}}$} & \textbf{Pol.\ Viol.} & \textbf{Overrides} \\
\hline
C1: Full        & .878$\pm$.007 & .931$\pm$.004 & .430$\pm$.039 & .683$\pm$.020 & .040$\pm$.018 & .151$\pm$.085 & 20.3$\pm$4.3 & 69.4$\pm$9.0 \\
C2: $-$Fairness & .885$\pm$.006 & .936$\pm$.003 & .368$\pm$.034 & .660$\pm$.018 & .019$\pm$.012 & .119$\pm$.065 & 20.3$\pm$4.3 & 0.0$\pm$0.0 \\
C3: $-$Policy   & .889$\pm$.009 & .939$\pm$.005 & .293$\pm$.034 & .630$\pm$.018 & .068$\pm$.016 & .247$\pm$.081 & 84.3$\pm$7.3 & 98.6$\pm$12.7 \\
C4: RF Only     & .902$\pm$.006 & .946$\pm$.004 & .357$\pm$.020 & .665$\pm$.011 & .062$\pm$.012 & .244$\pm$.078 & 78.0$\pm$9.0 & 0.0$\pm$0.0 \\
\hline
\end{tabular}
\end{table}

\paragraph{Effect of the Fairness Agent (C1 vs.\ C2).}
Disabling the counterfactual fairness agent removes the override rule that
reroutes gender-sensitive decisions. The most pronounced effect is on
specificity, which drops from $0.4300$ (C1) to $0.3681$ (C2), and balanced
accuracy, which decreases from $0.6832$ to $0.6603$. The override rule in
C1 triggers on $69.4 \pm 9.0$ applications per test split ($4.63\%$ of the
test set). Because the override rule can both approve and reject, it shifts
marginal decisions toward the correct class for defaulters, thereby raising
specificity. $\Delta_{\text{DP}}$ is higher in C1 ($0.0400$) than in C2
($0.0188$), and $\Delta_{\text{EO}}$ is likewise higher ($0.1506$ vs.\
$0.1193$). The counterfactual override thus increases group-level disparity
on both metrics. This is an empirical finding: the deterministic override
rule does not uniformly reduce demographic disparity.

\paragraph{Effect of the Policy Agent (C1 vs.\ C3).}
Removing the policy agent produces the largest change in regulatory
compliance. Policy violations rise from $20.3 \pm 4.3$ (C1) to
$84.3 \pm 7.3$ (C3), a fourfold increase. The approval rate correspondingly
increases from $89.4\%$ to $93.7\%$, as applications that would have been
blocked by hard regulatory rules (e.g., excessive debt-to-income ratios)
are now approved. Specificity drops from $0.4300$ to $0.2926$, confirming
that the policy agent is the primary mechanism by which the framework
rejects high-risk applicants. $\Delta_{\text{DP}}$ rises from $0.0400$ to
$0.0677$, and $\Delta_{\text{EO}}$ from $0.1506$ to $0.2466$, indicating
that the policy agent also contributes to fairness by uniformly blocking
non-compliant applicants regardless of demographic group.

\paragraph{Comparison with the RF Baseline (C1 vs.\ C4).}
The unconstrained Random Forest (C4) achieves the highest F1 ($0.9462$) and
accuracy ($0.9021$) among all configurations, but incurs $78.0 \pm 9.0$
policy violations per split. Specificity is $0.3569$ for C4 versus $0.4300$
for C1. $\Delta_{\text{DP}}$ is $0.0619$ (C4) versus $0.0400$ (C1), and
$\Delta_{\text{EO}}$ is $0.2439$ (C4) versus $0.1506$ (C1). The framework
thus reduces policy violations by $74.0\%$ and $\Delta_{\text{EO}}$ by
$38.3\%$ at the cost of a $1.6\%$ decrease in F1.

\subsection{Statistical Analysis}

Paired-sample $t$-tests (paired by seed, $n{=}10$) were conducted to assess
whether the observed differences between configurations are attributable to
chance variation in the train--test split. Table~\ref{tab:paired} reports
the results.

\begin{table}[htbp]
\centering
\caption{Paired-sample $t$-tests across 10 seeds. CI denotes the 95\%
confidence interval for the mean difference.}
\label{tab:paired}
\setlength{\tabcolsep}{3pt}
\footnotesize
\begin{tabular}{llrrr}
\hline
\textbf{Contrast} & \textbf{Metric} & \textbf{Mean Diff} & \textbf{95\% CI} & \textbf{$p$-value} \\
\hline
C1 vs C2 & $\Delta_{\text{DP}}$    & $+0.0211$ & $[0.0146,\; 0.0276]$ & $4.24 \times 10^{-5}$ \\
C1 vs C2 & $\Delta_{\text{EO}}$    & $+0.0313$ & $[0.0070,\; 0.0555]$ & $1.73 \times 10^{-2}$ \\
C1 vs C2 & Specificity             & $+0.0619$ & $[0.0517,\; 0.0721]$ & $2.49 \times 10^{-7}$ \\
C1 vs C2 & Balanced Accuracy       & $+0.0230$ & $[0.0177,\; 0.0282]$ & $3.92 \times 10^{-6}$ \\
\hline
C1 vs C3 & Policy Violations       & $-64.0$   & $[-67.50,\; -60.50]$ & $1.42 \times 10^{-11}$ \\
C1 vs C3 & Approval Rate           & $-0.0429$ & $[-0.0453,\; -0.0405]$ & $1.85 \times 10^{-11}$ \\
\hline
C1 vs C4 & Specificity             & $+0.0731$ & $[0.0454,\; 0.1008]$ & $2.09 \times 10^{-4}$ \\
C1 vs C4 & Balanced Accuracy       & $+0.0181$ & $[0.0044,\; 0.0317]$ & $1.53 \times 10^{-2}$ \\
C1 vs C4 & MCC                     & $-0.0486$ & $[-0.0732,\; -0.0241]$ & $1.52 \times 10^{-3}$ \\
C1 vs C4 & $\Delta_{\text{DP}}$    & $-0.0219$ & $[-0.0392,\; -0.0046]$ & $1.85 \times 10^{-2}$ \\
C1 vs C4 & $\Delta_{\text{EO}}$    & $-0.0934$ & $[-0.1568,\; -0.0299]$ & $8.81 \times 10^{-3}$ \\
\hline
\end{tabular}
\end{table}

Table~\ref{tab:paired} confirms that the specificity gain of C1 over C2
($+0.0619$, $p = 2.49 \times 10^{-7}$) and over C4 ($+0.0731$,
$p = 2.09 \times 10^{-4}$) is statistically significant at $\alpha = 0.01$.
The balanced accuracy advantage of C1 over C2 ($+0.0230$,
$p = 3.92 \times 10^{-6}$) is likewise significant. The reduction in policy
violations when comparing C1 to C3 ($-64.0$, $p = 1.42 \times 10^{-11}$) is
the strongest effect, with a confidence interval that excludes zero by a
wide margin.

The MCC of C1 is statistically significantly lower than that of C4
($-0.0486$, $p = 1.52 \times 10^{-3}$), confirming the predictive cost of
enforcement. On the fairness metrics, C1 exhibits lower $\Delta_{\text{DP}}$
($-0.0219$, $p = 1.85 \times 10^{-2}$) and lower $\Delta_{\text{EO}}$
($-0.0934$, $p = 8.81 \times 10^{-3}$) than C4, both significant at
$\alpha = 0.05$. However, C1 shows \emph{higher} $\Delta_{\text{DP}}$ and
$\Delta_{\text{EO}}$ than C2 (both $p < 0.05$), indicating that the
counterfactual override rule increases group-level disparity relative to
the system without it. This counterintuitive result arises because the
override rule is a deterministic threshold mechanism ($|\Delta_{\text{CF}}|
> \varepsilon$) that can systematically alter approval rates for one
demographic group more than the other, depending on the distribution of
counterfactual score differences across groups.

\subsection{Threshold Sensitivity Analysis}

The counterfactual fairness threshold $\varepsilon$ was swept over
$\{5.0, 5.5, 6.0, 6.5, 7.0, 7.5, 8.0\}$ on a held-out \emph{validation}
split (10\% of total data), disjoint from both the training set (60\%) and
the test set (30\%). The Random Forest was retrained on the 60\% training
portion for each of 10 seeds. Table~\ref{tab:sweep} reports the results;
Table~\ref{tab:trend} reports paired tests between the extreme values
$\varepsilon = 5.0$ and $\varepsilon = 8.0$.

\begin{table}[htbp]
\centering
\caption{Threshold sensitivity on the validation split (mean $\pm$ std,
$n{=}10$ seeds). Policy violations are constant across $\varepsilon$ because
the policy agent operates independently of the fairness threshold.}
\label{tab:sweep}
\setlength{\tabcolsep}{2.5pt}
\footnotesize
\begin{tabular}{ccccccc}
\hline
$\varepsilon$ & \textbf{F1} & \textbf{Bal.\ Acc.} & \textbf{Spec.} & \textbf{$\Delta_{\text{DP}}$} & \textbf{$\Delta_{\text{EO}}$} & \textbf{Override Rate} \\
\hline
5.0 & $.935 \pm .007$ & $.722 \pm .039$ & $.510 \pm .077$ & $.059 \pm .031$ & $.146 \pm .111$ & $5.48\%$ \\
5.5 & $.936 \pm .007$ & $.720 \pm .041$ & $.503 \pm .082$ & $.058 \pm .031$ & $.145 \pm .111$ & $4.80\%$ \\
6.0 & $.937 \pm .007$ & $.720 \pm .041$ & $.502 \pm .081$ & $.058 \pm .030$ & $.149 \pm .108$ & $4.40\%$ \\
6.5 & $.938 \pm .007$ & $.719 \pm .041$ & $.498 \pm .083$ & $.055 \pm .030$ & $.142 \pm .107$ & $3.98\%$ \\
7.0 & $.938 \pm .007$ & $.718 \pm .040$ & $.494 \pm .079$ & $.053 \pm .029$ & $.150 \pm .106$ & $3.68\%$ \\
7.5 & $.939 \pm .007$ & $.716 \pm .038$ & $.488 \pm .076$ & $.049 \pm .027$ & $.137 \pm .102$ & $3.20\%$ \\
8.0 & $.940 \pm .008$ & $.715 \pm .040$ & $.484 \pm .078$ & $.046 \pm .028$ & $.130 \pm .106$ & $2.70\%$ \\
\hline
\end{tabular}
\end{table}

\begin{table}[htbp]
\centering
\caption{Paired tests between $\varepsilon = 5.0$ and $\varepsilon = 8.0$
($n{=}10$ seeds), with Spearman rank correlations between $\varepsilon$ and
each metric across all 70 seed--threshold observations.}
\label{tab:trend}
\footnotesize
\begin{tabular}{lrrrcc}
\hline
\textbf{Metric} & \textbf{Mean Diff} & \textbf{95\% CI} & \textbf{$p_{\text{paired}}$} & \textbf{$\rho_s$} & \textbf{$p_{\text{Spearman}}$} \\
\hline
$\Delta_{\text{DP}}$  & $+0.0128$ & $[0.0049,\; 0.0207]$ & $5.24 \times 10^{-3}$ & $-0.191$ & $0.112$ \\
$\Delta_{\text{EO}}$  & $+0.0169$ & $[-0.0276,\; 0.0613]$ & $0.413$ & $-0.055$ & $0.654$ \\
Specificity           & $+0.0256$ & $[0.0114,\; 0.0398]$ & $2.73 \times 10^{-3}$ & $-0.183$ & $0.129$ \\
Balanced Accuracy     & $+0.0075$ & $[0.0002,\; 0.0148]$ & $4.52 \times 10^{-2}$ & $-0.065$ & $0.594$ \\
\hline
\end{tabular}
\end{table}

Table~\ref{tab:sweep} shows that increasing $\varepsilon$ from 5.0 to 8.0
monotonically decreases the override trigger rate from $5.48\%$ to $2.70\%$,
as fewer counterfactual score differences exceed the threshold. This
reduction in interventions is accompanied by a decrease in specificity
($0.510 \to 0.484$) and balanced accuracy ($0.722 \to 0.715$), and
a decrease in $\Delta_{\text{DP}}$ ($0.059 \to 0.046$). F1 increases
marginally ($0.935 \to 0.940$) as fewer approved applications are
overridden. Policy violations remain constant at $6.4 \pm 2.5$ across all
thresholds, confirming that the policy agent operates independently of
$\varepsilon$.

Table~\ref{tab:trend} quantifies these trends. The paired difference in
$\Delta_{\text{DP}}$ between $\varepsilon = 5.0$ and $\varepsilon = 8.0$ is
$+0.0128$ ($p = 5.24 \times 10^{-3}$), indicating that a lower threshold
produces statistically significantly higher demographic parity disparity.
Specificity likewise decreases significantly when the threshold is raised
($+0.0256$, $p = 2.73 \times 10^{-3}$). The balanced accuracy difference is
marginally significant ($p = 0.045$). The $\Delta_{\text{EO}}$ difference is
not statistically significant ($p = 0.413$). Spearman correlations between
$\varepsilon$ and each metric across all 70 observations are weak
($|\rho_s| < 0.20$) and none reach significance at $\alpha = 0.05$,
reflecting the high inter-seed variance relative to the monotonic trend.

The sweep confirms that $\varepsilon$ governs a continuous trade-off:
lower thresholds increase the frequency of override interventions, which
raises specificity and balanced accuracy but also increases
$\Delta_{\text{DP}}$. The current operating point ($\varepsilon = 6.0$)
represents a selection within this continuum; the data do not identify a
single threshold that simultaneously minimises all objectives.

\subsection{Discussion}

The experimental results expose a fundamental tension between predictive
accuracy and regulatory governance in automated lending. The unconstrained
Random Forest baseline achieves the highest F1 score and accuracy among all
configurations, and even a trivial always-approve predictor surpasses the
proposed framework on both metrics due to class imbalance. These findings
underscore that accuracy and F1 alone are insufficient evaluation criteria
for credit-risk systems operating under regulatory constraints.

The ablation study reveals that the policy agent contributes the largest
single effect: its removal causes policy violations to increase from $20.3$
to $84.3$ ($p = 1.42 \times 10^{-11}$) and specificity to drop from
$0.4300$ to $0.2926$. The policy agent also indirectly reduces
$\Delta_{\text{DP}}$ and $\Delta_{\text{EO}}$, because hard regulatory rules
(e.g., debt-to-income ceilings) reject non-compliant applicants uniformly
across demographic groups. The fairness agent contributes a statistically
significant gain in specificity ($+0.0619$, $p = 2.49 \times 10^{-7}$) and
balanced accuracy ($+0.0230$, $p = 3.92 \times 10^{-6}$) but, contrary to
its design intent, increases $\Delta_{\text{DP}}$ by $0.0211$
($p = 4.24 \times 10^{-5}$) and $\Delta_{\text{EO}}$ by $0.0313$
($p = 1.73 \times 10^{-2}$). The counterfactual override is a deterministic
threshold rule that applies the same $\varepsilon$ to all applicants; when
the distribution of counterfactual score differences differs between
demographic groups, the override can disproportionately alter one group's
approval rate.

The threshold sensitivity analysis confirms that the fairness--specificity
trade-off is monotonic and continuous. Lowering $\varepsilon$ increases both
the frequency of override interventions and the resulting specificity, but
also raises $\Delta_{\text{DP}}$. The absence of a Pareto-dominant operating
point means that the choice of $\varepsilon$ is inherently a policy
decision, requiring domain-specific balancing of competing objectives.

\paragraph{Limitations.}
Several limitations qualify these findings.
(i)~The dataset is synthetic and generated from a parametric causal model;
the distributional assumptions (e.g., log-normal income, linear logit for
default probability) may not reflect real-world credit populations.
(ii)~The counterfactual fairness audit considers only a single binary
protected attribute (gender). Intersectional fairness across multiple
attributes remains unexamined.
(iii)~The override rule is deterministic---it applies a fixed threshold
$\varepsilon$ without adapting to the applicant's position in the score
distribution. An adaptive or calibrated mechanism might avoid the observed
increase in $\Delta_{\text{DP}}$.
(iv)~The evaluation uses a Random Forest with fixed hyperparameters
($n_{\text{estimators}} = 100$); the relative performance of the agentic
layers may differ under alternative base classifiers or with hyperparameter
tuning.
(v)~The multi-agent orchestration is evaluated offline; latency, failure
modes, and integration costs in a production deployment are not assessed.
